\documentclass[11pt]{article}
\usepackage[margin=1in]{geometry}
\usepackage{xcolor}
\usepackage[breakable,listings]{tcolorbox}
\tcbuselibrary{listings}
\usepackage{hyperref}
\usepackage{enumitem}
\usepackage{listings}
\usepackage{amsmath}
\usepackage{booktabs}
\usepackage{array}

% Define colors
\definecolor{osaorange}{RGB}{220,115,38}

% Configure hyperref
\hypersetup{
    colorlinks=true,
    linkcolor=blue,
    urlcolor=blue
}

% Code environment
\newtcblisting{rcode}{
  colback=white,
  colframe=gray!50,
  boxrule=0.5pt,
  arc=0pt,
  left=3pt,
  right=3pt,
  top=3pt,
  bottom=3pt,
  width=\textwidth,
  breakable,
  listing only,
  listing options={
    basicstyle=\small\ttfamily,
    breaklines=true,
    breakatwhitespace=true,
    columns=flexible,
    keepspaces=true,
    showstringspaces=false
  }
}

\title{}
\author{}
\date{}

\begin{document}

% Header box
\begin{tcolorbox}[colback=osaorange,colframe=black,boxrule=2pt,arc=0pt,left=3pt,right=3pt,top=6pt,bottom=6pt,boxsep=2pt]
\centering
\textcolor{white}{\textbf{\Large H524 Introduction to Biostatistics}}\\
\textcolor{white}{\textbf{\Large Week 6: Power, Two-Sample Tests, and ANOVA}}\\
\textcolor{white}{\textbf{\Large Homework Assignment}}\\
\textcolor{white}{\textbf{\Large Fall 2025}}
\end{tcolorbox}

\vspace{8pt}

\noindent\textbf{Due Date:} End of Week 6\\
\textbf{Total Points:} 28 points\\
\textbf{Format:} Submit through Canvas\\
\textbf{Time Limit:} 60 minutes

\section*{Instructions}

\begin{itemize}
\item Show all work for full credit on numerical questions
\item Round numerical answers to 2 decimal places unless otherwise specified
\item Include R code where specified
\item You may use R to perform calculations
\item For multiple choice, select the BEST answer
\end{itemize}

\section*{Multiple Choice Questions (10 points total, 2 points each)}

\textbf{Question 1:} Statistical power is defined as:

\begin{enumerate}[label=\Alph*.]
\item The probability of making a Type I error
\item The probability of correctly rejecting a false null hypothesis
\item The probability of failing to reject a true null hypothesis
\item The significance level of the test
\end{enumerate}

\vspace{0.5cm}

\textbf{Question 2:} Which of the following will \textit{increase} the statistical power of a study?

\begin{enumerate}[label=\Alph*.]
\item Decreasing the sample size
\item Increasing the population variance
\item Increasing the significance level (alpha)
\item Decreasing the effect size
\end{enumerate}

\vspace{0.5cm}

\textbf{Question 3:} A researcher measures blood pressure in 15 patients before and after a medication intervention. Which test should be used?

\begin{enumerate}[label=\Alph*.]
\item One-sample t-test
\item Independent two-sample t-test
\item Paired t-test
\item One-way ANOVA
\end{enumerate}

\vspace{0.5cm}

\textbf{Question 4:} When comparing means across 4 independent groups, why is ANOVA preferred over multiple pairwise t-tests?

\begin{enumerate}[label=\Alph*.]
\item ANOVA is easier to calculate
\item ANOVA requires smaller sample sizes
\item ANOVA controls the family-wise Type I error rate
\item ANOVA doesn't require assumptions
\end{enumerate}

\vspace{0.5cm}

\textbf{Question 5:} An ANOVA test comparing 5 groups yields F = 3.45 with p = 0.02. What can you conclude?

\begin{enumerate}[label=\Alph*.]
\item All five group means are significantly different from each other
\item At least one group mean differs from the others
\item The first group differs significantly from the last group
\item Groups 1 and 2 have equal means
\end{enumerate}

\newpage

\section*{Numerical Questions (15 points total)}

\subsection*{Question 6: Sample Size for Power (3 points)}

A researcher is planning a study to test whether a new medication reduces systolic blood pressure. Based on pilot data, the population standard deviation is estimated to be $\sigma = 20$ mmHg. The researcher wants to detect a mean reduction of 10 mmHg from the population mean of 145 mmHg.

\vspace{0.3cm}

Using a significance level of $\alpha = 0.05$ (two-sided) and desired power of 0.80, calculate the required sample size using R's \texttt{power.t.test()} function for a one-sample t-test.

\vspace{0.3cm}

\textbf{Show your R code:}

\vspace{3cm}

\textbf{What is the required sample size (round UP to the nearest whole number)?}

\vspace{0.3cm}

\textbf{Answer:} \underline{\hspace{3cm}} subjects

\vspace{1cm}

\subsection*{Question 7: Two-Sample t-Test and ANOVA (3 points)}

A researcher compares the effectiveness of three different exercise programs on weight loss (in kg) over 8 weeks. The data are:

\vspace{0.3cm}

\begin{center}
\begin{tabular}{lll}
\toprule
Program A & Program B & Program C \\
\midrule
4.2 & 6.1 & 7.8 \\
3.8 & 5.8 & 8.2 \\
4.5 & 6.4 & 7.5 \\
4.1 & 5.9 & 8.0 \\
4.3 & 6.2 & 7.9 \\
\bottomrule
\end{tabular}
\end{center}

\vspace{0.3cm}

\textbf{Part A (1.5 points):} Perform a one-way ANOVA to test if there is a significant difference in mean weight loss across the three programs. Use $\alpha = 0.05$.

\vspace{0.3cm}

\textbf{Show your R code to create the data and perform ANOVA:}

\vspace{5cm}

\textbf{What is the F-statistic (round to 2 decimal places)?}

\vspace{0.3cm}

\textbf{Answer:} F = \underline{\hspace{3cm}}

\vspace{0.5cm}

\textbf{Part B (1.5 points):} Compare ONLY Program A vs. Program C using a two-sample t-test (assume equal variances, two-sided test).

\vspace{0.3cm}

\textbf{Show your R code:}

\vspace{4cm}

\textbf{What is the p-value (round to 4 decimal places)?}

\vspace{0.3cm}

\textbf{Answer:} p-value = \underline{\hspace{3cm}}

\vspace{1cm}

\subsection*{Question 8: Pooled Standard Deviation (3 points)}

Two independent samples are collected to compare cholesterol levels between two diets:

\begin{itemize}
\item Diet 1: $n_1 = 12$, $s_1 = 18.5$ mg/dL
\item Diet 2: $n_2 = 15$, $s_2 = 21.3$ mg/dL
\end{itemize}

Calculate the pooled standard deviation ($s_p$) using the formula:

$$s_p = \sqrt{\frac{(n_1-1)s_1^2 + (n_2-1)s_2^2}{n_1 + n_2 - 2}}$$

\vspace{0.3cm}

\textbf{Show your calculation (R code or by hand):}

\vspace{3cm}

\textbf{Answer:} $s_p$ = \underline{\hspace{3cm}} mg/dL (round to 2 decimal places)

\vspace{1cm}

\subsection*{Question 9: Test Statistic Calculation (3 points)}

A researcher compares mean test scores between two teaching methods. The summary statistics are:

\begin{itemize}
\item Method 1: $n_1 = 20$, $\bar{x}_1 = 85.3$, $s_1 = 12.4$
\item Method 2: $n_2 = 18$, $\bar{x}_2 = 78.6$, $s_2 = 14.1$
\end{itemize}

Calculate the test statistic (t-value) for a two-sample t-test (assume equal variances) using the formula:

$$t = \frac{\bar{x}_1 - \bar{x}_2}{s_p \sqrt{\frac{1}{n_1} + \frac{1}{n_2}}}$$

where $s_p = \sqrt{\frac{(n_1-1)s_1^2 + (n_2-1)s_2^2}{n_1 + n_2 - 2}}$

\vspace{0.3cm}

\textbf{Show your calculation (R code or by hand):}

\vspace{4cm}

\textbf{Answer:} t = \underline{\hspace{3cm}} (round to 2 decimal places)

\vspace{1cm}

\subsection*{Question 10: Confidence Interval for Mean Difference (3 points)}

Using the data from Question 9 (Teaching Methods), construct a 95\% confidence interval for the mean difference in test scores ($\mu_1 - \mu_2$).

Use the formula:

$$\text{CI} = (\bar{x}_1 - \bar{x}_2) \pm t_{\alpha/2, df} \times SE$$

where $SE = s_p \sqrt{\frac{1}{n_1} + \frac{1}{n_2}}$ and $df = n_1 + n_2 - 2$

\textbf{Or} use R to calculate (show your work either way):

\vspace{4cm}

\textbf{Answer:} 95\% CI = (\underline{\hspace{2cm}}, \underline{\hspace{2cm}}) (round to 2 decimal places)

\vspace{0.5cm}

\textbf{Interpretation:} Based on this confidence interval, is there evidence of a significant difference between the two teaching methods at $\alpha = 0.05$? Explain briefly.

\vspace{2cm}

\subsection*{Question 11: Power Calculation (3 points)}

A researcher is planning a one-sample study to test whether a new intervention changes blood glucose levels. Based on prior research:

\begin{itemize}
\item Expected sample size: $n = 30$
\item Expected mean difference to detect: $\delta = 12$ mg/dL
\item Population standard deviation: $\sigma = 18$ mg/dL
\item Significance level: $\alpha = 0.05$ (two-sided)
\end{itemize}

Calculate the statistical power of this study using R's \texttt{power.t.test()} function.

\vspace{0.3cm}

\textbf{Show your R code:}

\vspace{4cm}

\textbf{Answer:} Power = \underline{\hspace{3cm}} (round to 2 decimal places)

\vspace{0.5cm}

\textbf{Note:} Express your answer as a proportion (e.g., 0.85), not as a percentage.

\vspace{1cm}

\end{document}
